% Table 1 for the proceeding. Every number traceable to MAGI_state_reference.pdf.
% R = (MAGI counts / MAGI injected crossings) / (real counts / real shot particles),
% so R = 1 means a MAGI crossing produces detector events at exactly the rate a
% real recorded crossing does. eps = R/chi is the efficiency factor; a perfect
% surrogate returns eps = 1/chi, not 1.
\begin{table*}[t]
\centering
\caption{MAGI v0.8.2 against full Geant4, at the detector, in the same mass
model. Counts are events with summed deposit in the quoted band. Uncertainties
are Poisson on both arms. The DM1.2 0.1--30\,keV row is an integrated rate, not
a spectral comparison: DM127 is an anti-coincidence detector and the reference
holds only 41 events there.}
\label{tab:validation}
\begin{tabular}{llrrrl}
\hline\hline
Source / setup & Band & ref.\ counts & MAGI counts & $R$ & notes \\
\hline
\multicolumn{6}{l}{\textit{DM1.2 laboratory cryostat} --- $2\times10^{8}\,\mu$ (iso) vs $1.2\times10^{6}$ injected crossings, 3-seed ensemble} \\
 & all deposits      & 1733 & 2019 & $0.970 \pm 0.032$ & \\
 & 100--300\,keV (MIP) & 1325 & 1589 & $\mathbf{0.999 \pm 0.037}$ & headline; $\varepsilon = 199.9 \pm 7.4$ vs ideal 200.13 \\
 & 300--1000\,keV    &  310 &  323 & $0.868 \pm 0.069$ & \\
 & 30--100\,keV      &   42 &   46 & $0.912 \pm 0.195$ & \\
 & 0.1--30\,keV      &   41 &   43 & $0.873 \pm 0.191$ & reference-limited ($\pm15.6\,\%$) \\
 & $>1000$\,keV      &   13 &   17 & $1.089 \pm 0.401$ & \\
\hline
\multicolumn{6}{l}{\textit{SRON X-IFU model, cosmic rays}} \\
 & 0.1--30\,keV, normalization-free &   51 &  500 & $0.848 \pm 0.125$ & real crossings vs MAGI, same geometry \\
 & 1--20\,keV, Det.~1 flux          & \multicolumn{2}{c}{---} & $0.76$ & $\varepsilon = 124.8 \pm 5.6$ vs ideal 164.4 \\
 & 1--20\,keV, mean over 6 detectors& \multicolumn{2}{c}{---} & $0.79$ & $\varepsilon = 124.1 \pm 15.2$ \\
\hline
\multicolumn{6}{l}{\textit{$^{40}$K, SRON BuildingModel} --- $2\times10^{8}$ injected on both arms;
\textbf{method validation only, see note}} \\
 & 0.1--30\,keV      &   17 &   19 & $1.118 \pm 0.373$ & KS $D=0.238$, $p=0.60$ \\
\hline\hline
\end{tabular}

\vspace{2mm}
{\footnotesize \textbf{Note on $^{40}$K.} Both arms of that row draw on the same
training population, which was recorded through a surface-crossing test whose
centre was displaced by 2.5\,m and therefore retains only back-scattered photons
($100.0\,\%$ ingoing w.r.t.\ the wrong centre; median $211$\,keV against a
$217.5$\,keV Compton back-scatter energy). The row is a valid closure test of
MAGI against its own training distribution and is \emph{not} a $^{40}$K
background result. $^{226}$Ra and $^{232}$Th share the defect. Regeneration is
scheduled for September 2026.}

\caption{Crossing-level fidelity and cost, cosmic rays on the SRON model.}
\label{tab:fidelity}
\begin{tabular}{lll}
\hline\hline
Quantity & Value & Comment \\
\hline
Correlation residual, mean $|\Delta\rho|$ & 0.015 (max 0.042) & 10 feature pairs, Fig.~2 \\
Species fractions & correct to $0.3\,\%$ & $\gamma$, $e^{-}$, $e^{+}$, $\mu^{-}$ \\
Energy marginals, per decade & within a few \% & 0.1\,MeV--10\,GeV \\
Geometry marginals, RMS per bin & 5--11\,\% & vs a $0.44\,\%$ Poisson floor \\
Spectral shape, SRON Det.~1 & KS $p = 0.20$ & indistinguishable from reference \\
\hline
Statistical amplification (DM1.2) & $194.2\times$ & ideal $1/\chi = 200.13$ \\
Transport cost ratio & $95.8\times$ & crossing vs thrown muon \\
Net throughput gain & $2.03\times$ & measured, same machine \\
Outer / inner cost split & $52\,\% / 48\,\%$ & MAGI eliminates the outer \\
Break-even, design campaign & 2nd variant & one training set, $N$ inner designs \\
$^{40}$K sample amplification & $287\times$ & $1.87\times10^{6} \to 5.38\times10^{8}$ \\
\hline\hline
\end{tabular}
\end{table*}
